\chapter{Flujos}
\label{cap:flujos}

Tres recorridos completos, elegidos porque cada uno atraviesa una frontera
distinta: el modelo de visión, el protocolo MCP y la concurrencia sobre una fila
de la base de datos. Los tres diagramas se generan con PlantUML desde
\pkpath{diagramas/src/}.

\section{Escaneo de una carta}

\begin{figure}[p]
  \centering
  \includegraphics[width=\linewidth,height=0.84\textheight,keepaspectratio]{seq-escaneo.pdf}
  \caption{\pkcode{POST /api/v1/scans} y la confirmación posterior.}
  \label{fig:seq-escaneo}
\end{figure}

Tres decisiones separadas, y sólo la primera es del modelo:

\begin{enumerate}[leftmargin=1.4em,itemsep=3pt,topsep=3pt]
  \item \textbf{El modelo lee lo impreso.} El \emph{prompt} de visión dice
        explícitamente «no infieras ni identifiques de qué carta se trata; otro
        sistema hace eso». Devuelve nombre, número, total del set y HP.
  \item \textbf{El código decide qué carta es.} \pkpath{services/resolve.py}
        puntúa cada señal (figura~\ref{fig:resolve}) en lugar de filtrar por
        ninguna, y clasifica el resultado en \pkcode{resolved},
        \pkcode{ambiguous} o \pkcode{failed}.
  \item \textbf{La persona decide qué se guarda}, y sólo cuando la ambigüedad es
        real.
\end{enumerate}

Dos detalles operativos: la imagen se guarda \emph{antes} de llamar al modelo,
para que un fallo de visión deje algo que reintentar; y si no hay modelo
configurado, \pkcode{create\_scan} captura el error y sigue con un
\pkcode{CardReading} vacío, de modo que el escaneo degrada a entrada manual en
vez de fallar.

\section{Un turno del asistente sobre MCP}

\begin{figure}[p]
  \centering
  \includegraphics[width=\linewidth,height=0.84\textheight,keepaspectratio]{seq-chat-mcp.pdf}
  \caption{\pkcode{POST /api/v1/chat}: un turno completo, con sesión MCP propia y
    respuesta en streaming.}
  \label{fig:seq-chat}
\end{figure}

Lo importante del diagrama es que la aplicación se llama a sí misma por el
protocolo público. El chat no consulta \pkpath{services/} directamente: abre una
sesión MCP contra su propio servidor, con \emph{el token del usuario}, y de ahí
en adelante es un cliente MCP más. Las consecuencias:

\begin{itemize}[leftmargin=1.1em,itemsep=3pt,topsep=3pt]
  \item El agente no puede hacer nada que un asistente externo no pueda hacer.
        No hay una ruta privilegiada.
  \item El agente actúa \emph{como el llamante}, no como una identidad de
        servicio: el aislamiento por usuario lo aplica la misma verificación de
        token que protege el REST.
  \item Un turno está acotado: ocho llamadas a herramienta y diez peticiones al
        modelo. «Un turno que no termina nunca es peor que una respuesta
        incompleta.»
  \item El historial se recorta a veinte turnos y se descartan los viejos en
        lugar de resumirlos, porque las herramientas releen datos vivos en cada
        turno.
\end{itemize}

\section{Tomar una publicación del tablón}

\begin{figure}[p]
  \centering
  \includegraphics[width=\linewidth,height=0.84\textheight,keepaspectratio]{seq-tablon.pdf}
  \caption{\pkcode{POST /api/v1/trades/listings/\{id\}/accept}: la reserva
    atómica y el trato que queda escrito.}
  \label{fig:seq-tablon}
\end{figure}

Este es el único punto del sistema con competencia real por un recurso: una
publicación la puede tomar cualquiera, y sólo una vez. La reserva es una sola
sentencia:

\begin{lstlisting}[language=Python]
claimed = (await db.execute(
    update(TradeListing)
    .where(TradeListing.id == listing.id,
           TradeListing.status == ListingStatus.OPEN)
    .values(status=ListingStatus.TAKEN, taken_at=func.now())
    .returning(TradeListing.id)
)).scalar_one_or_none()

if claimed is None:
    raise ListingError("This listing was already taken")
\end{lstlisting}

La condición sobre el estado dentro del \pkcode{WHERE} es el cerrojo: dos
peticiones simultáneas producen un ganador —el que ve una fila en el
\pkcode{RETURNING}— y la otra recibe 409. No hace falta bloqueo explícito ni un
nivel de aislamiento especial.

Antes de la reserva, ambos lados se revalidan contra los sobrantes de
\emph{ahora}, no contra los de cuando se publicó. Después, se escribe una
\pkcode{trade\_offer} ya aceptada y la publicación queda enlazada a ella, con lo
que un trato del tablón y uno negociado son la misma cosa desde ese momento.

\begin{pkdecision}[Aceptar no mueve cartas]
Ni aquí ni al aceptar una oferta cambia una sola fila de \pkcode{collection\_item}.
La colección es lo que alguien tiene físicamente, y sólo quien tiene la carta
sabe si el intercambio ocurrió. Es exactamente la misma regla que impide al
asistente archivar cartas.
\end{pkdecision}
